\chapter{Endliche Kanonisierung und homomorphe Bilder}

\section{Multiplikationstafeln und endliche Entscheidung}

\begin{definition}[Nummerierte Multiplikationstafel]
Für $n\geq0$ sei $[n]=\{0,\ldots,n-1\}$, mit $[0]=\varnothing$.
Eine \emph{Tafel} ist eine Abbildung $m:[n]^2\to[n]$.
Sie heißt assoziativ, wenn
\[
 m(m(i,j),k)=m(i,m(j,k))\qquad(i,j,k\in[n]).
\]
Die Menge der assoziativen Tafeln sei $\mathcal A_n$; mit
$\mathfrak S_n$ bezeichnen wir die Bijektionen von $[n]$ auf sich selbst.
Für $m\in\mathcal A_n$ und $\sigma\in\mathfrak S_n$ setzen wir
\[
 (\sigma\cdot m)(i,j)
 =\sigma\bigl(m(\sigma^{-1}(i),\sigma^{-1}(j))\bigr).
\]
Zu einer Bijektion $b:[n]\to S$ gehört die nummerierte Tafel
\[
 m_b(i,j)=b^{-1}\bigl(b(i)b(j)\bigr)\qquad(i,j\in[n]).
\]
\end{definition}

\begin{definition}[Kanonischer Tafelcode]
Es sei $\operatorname{vec}(m)$ das zeilenweise gelesene Tupel der $n^2$
Tafeleinträge. Der \emph{kanonische Code} ist
\[
 c(m)=\min_{\mathrm{lex}}\{
       \operatorname{vec}(\sigma\cdot m):\sigma\in\mathfrak S_n\}.
\]
Für eine beliebige endliche Halbgruppe $S$ sei
$c(S)=(|S|,c(m_b))$, wobei $b:[|S|]\to S$ eine beliebige Bijektion ist.
\end{definition}

\begin{theorem}[Endliche Kanonisierung und Entscheidungsverfahren]
Für jede endliche Halbgruppe $S$ mit $n=|S|$ existiert das
lexikographische Minimum der Tupel ihrer nummerierten Tafeln.
Der kanonische Code $c(S)$ lässt sich durch Prüfung von höchstens
$n!$ Umbenennungen berechnen.

Die Isomorphie zweier endlicher Halbgruppen ist entscheidbar;
bei gleicher Größe $n$ genügt die Prüfung von höchstens $n!$
Bijektionen auf Erhaltung sämtlicher Produkte.
Die Assoziativität einer Tafel auf $[n]$ ist durch höchstens
$n^3$ Tripeltests entscheidbar. Für jede feste Größe lassen sich
durch endliche Aufzählung und Isomorphieprüfung sämtliche
Isomorphieklassen vollständig erfassen.
\cite{projecttables}
\end{theorem}

\section{Hebungen aus homomorphen Bildern}

\begin{theorem}[Hebung von Idempotenten und Untergruppen]
Für jeden surjektiven Homomorphismus $f:S\to T$ endlicher Halbgruppen
und jedes $e\in E(T)$ gibt es ein $d\in E(S)$ mit
\[
 f(d)=e.
\]
Zu jeder Untergruppe $G\subseteq T$ gibt es eine Untergruppe
$H\subseteq S$ mit $f[H]=G$.
\cite{branco2020}
\end{theorem}

\begin{theorem}[Hebung regulärer Klassen nach Rhodes]
\label{thm:b37-guide-rhodes-lifting}
Für jeden surjektiven Homomorphismus $f:S\to T$ endlicher Halbgruppen
und jede reguläre $\mathcal J$-Klasse $J$ von $T$ gibt es unter den
$\mathcal J$-Klassen von $S$, die in $f^{-1}(J)$ liegen, eine bezüglich
der Klassenordnung minimale Klasse $\widetilde J$, die regulär ist.

Zu jeder maximalen Untergruppe $G$ von $T$, die in $J$ liegt,
gibt es eine maximale Untergruppe $H$ von $S$, die in
$\widetilde J$ liegt und $f[H]=G$ erfüllt.
\cite[Proposition~3.11]{AlmeidaMargolisSteinbergVolkov2009}
\end{theorem}

\clearpage
\section{Faktorisierung surjektiver Homomorphismen}

\begin{definition}[Homomorphismen und $\mathcal H$-Klassen]
Ein Homomorphismus $f:A\to B$ heißt \emph{$\mathcal H$-injektiv}, wenn
\[
 x\mathrel{\mathcal H_A}y,\quad f(x)=f(y)
 \quad\Longrightarrow\quad x=y
 \qquad(x,y\in A).
\]
Er heißt \emph{$\mathcal H$-Homomorphismus}, wenn jede seiner nichtleeren Fasern
in einer $\mathcal H$-Klasse von $A$ liegt, also
\[
 f(x)=f(y)\quad\Longrightarrow\quad x\mathrel{\mathcal H_A}y
 \qquad(x,y\in A).
\]
\end{definition}

\begin{theorem}[Homomorphismensatz von Rhodes]
\label{thm:b37-guide-rhodes-factorization}
Jeder surjektive Homomorphismus $f:S\to T$ endlicher Halbgruppen
besitzt eine Faktorisierung
\[
 S=S_0\xrightarrow{f_1}S_1\xrightarrow{f_2}\cdots
 \xrightarrow{f_m}S_m=T,\qquad
 f=f_m\circ\cdots\circ f_1
\]
mit $m\geq1$, endlichen Halbgruppen $S_i$ und surjektiven
Homomorphismen $f_i$, von denen jeder $\mathcal H$-injektiv
oder ein $\mathcal H$-Homomorphismus ist.
\cite{Rhodes1967Homomorphism,Rhodes1970Survey}
\end{theorem}

\section{Endliche Überlagerungen}

\begin{definition}[Rechter Stabilisator]
Für eine Halbgruppe $T$ und $a\in T$ heißt
\[
 \operatorname{Stab}_{R}^{T}(a)=\{u\in T:au=a\}
\]
der \emph{rechte Stabilisator} von $a$ in $T$.
Diese Menge darf leer sein.
\end{definition}

\begin{theorem}[Überlagerungssatz von Le Saëc--Pin--Weil]
\label{thm:b37-guide-stabilizer-cover}
Zu jeder endlichen Halbgruppe $S$ gibt es eine endliche Halbgruppe $T$
und einen surjektiven Homomorphismus $\pi:T\to S$, so daß für jedes
$a\in T$ und alle $u,v\in\operatorname{Stab}_{R}^{T}(a)$ gilt:
\[
 u^2=u,\qquad uvu=uv.
\]
\cite[Satz~1.1]{LeSaecPinWeil1991}
\end{theorem}

\clearpage
\section{Endliche Normalformen über einem vorgegebenen Alphabet}

\begin{definition}[Gewichtete Wörter]
Ein \emph{Alphabet} $A$ ist eine endliche Menge. Mit $A^*$ bezeichnen
wir die endlichen Wörter über $A$, einschließlich des leeren Wortes
$\varepsilon$, mit Konkatenation als Produkt.
Eine \emph{Gewichtung} $w:A\to\mathbb N_{>0}$ wird durch
\[
 w(a_1\cdots a_k)=\sum_{i=1}^k w(a_i),\qquad w(\varepsilon)=0
\]
auf Wörter fortgesetzt.
\end{definition}

\begin{definition}[Umschreibungsregeln und Konfluenz]
Ein \emph{Regelsystem} ist eine endliche Teilmenge $R\subseteq A^*\times A^*$.
Für $(\ell,r)\in R$ und $x,y\in A^*$ ist
$x\ell y\to_R xry$ ein \emph{Umschreibungsschritt}.
Die Relation $\to_R^*$ erlaubt endlich viele solche Schritte,
einschließlich null Schritten.
Das System heißt \emph{konfluent}, wenn zu
\[
 u\to_R^*v,\qquad u\to_R^*z
\]
stets ein $t\in A^*$ mit $v\to_R^*t$ und $z\to_R^*t$ existiert.
Es heißt \emph{gewichtssenkend}, wenn $w(\ell)>w(r)$ für alle
$(\ell,r)\in R$ gilt. Mit $\operatorname{IRR}_R(A^*)$ bezeichnen wir
die Wörter, auf die kein Umschreibungsschritt anwendbar ist.
\end{definition}

\begin{theorem}[Church--Rosser-Satz für endliche Halbgruppen]
\label{thm:b37-guide-church-rosser}
Für jede endliche Halbgruppe $S$, jedes endliche Alphabet $A$,
jede Gewichtung $w:A\to\mathbb N_{>0}$ und jeden Homomorphismus
$\varphi:A^*\to S^1$ gibt es ein endliches, konfluent
gewichtssenkendes Regelsystem $R\subseteq A^*\times A^*$ mit
\[
 |\operatorname{IRR}_R(A^*)|<\infty,
 \qquad
 \varphi(\ell)=\varphi(r)\quad((\ell,r)\in R).
\]
\cite[Satz~5.1]{DiekertKufleitnerReinhardtWalter2015}
\end{theorem}

